%%
% 附录：核心源代码
%%

% 配置 listings 样式
\lstset{
    basicstyle=\small\ttfamily,
    keywordstyle=\color{blue}\bfseries,
    commentstyle=\color{gray}\itshape,
    stringstyle=\color{red},
    numbers=left,
    numberstyle=\tiny\color{gray},
    numbersep=5pt,
    frame=single,
    breaklines=true,
    breakatwhitespace=true,
    showstringspaces=false,
    tabsize=4,
    captionpos=b,
    xleftmargin=2em,
    framexleftmargin=1.5em,
    literate={"}{{\textquotedbl}}1
             {'}{{\textquotesingle}}1,
    upquote=true,
    columns=flexible
}

\chapter{核心源代码}
\label{app:source-code}

\section{数据模型 (models.py)}
\label{sec:models}

\begin{lstlisting}[language=Python, caption={数据模型定义 (pathtracer/models.py)}]
"""PathTracer 数据模型: 定义执行路径追踪和覆盖率分析的核心数据结构"""

from dataclasses import dataclass, field
from typing import Set, Dict, List, Optional
from enum import Enum

# 枚举: 分支类型, 用于标识代码中的控制流结构
class BranchType(Enum):
    SEQUENTIAL = "sequential"  # 顺序执行
    IF = "if"                  # if 条件分支
    ELSE = "else"              # else 分支
    ELIF = "elif"              # elif 分支
    FOR = "for"                # for 循环
    WHILE = "while"            # while 循环
    TRY = "try"                # try 异常处理块
    EXCEPT = "except"          # except 异常捕获块

# 数据类: 代码分支点
# 记录单个分支的位置、类型和执行状态
@dataclass
class CodeBranch:
    file_path: str              # 源文件路径
    line_no: int                # 分支起始行号
    branch_type: BranchType     # 分支类型
    end_line: int = 0           # 分支结束行号
    is_executed: bool = False   # 是否被执行过

# 数据类: 函数调用记录
# 记录一次函数调用的完整信息, 支持嵌套调用树
@dataclass
class FunctionCall:
    function_name: str                    # 函数名
    file_path: str                        # 所在文件
    line_no: int                          # 调用行号
    call_depth: int                       # 调用深度 (0=顶层)
    arguments: Dict[str, str] = field(default_factory=dict)  # 参数及值
    return_value: Optional[str] = None    # 返回值
    children: List['FunctionCall'] = field(default_factory=list)  # 子调用

# 数据类: 执行路径
# 记录单个文件的完整执行轨迹
@dataclass
class ExecutionPath:
    file_path: str                                    # 源文件路径
    executed_lines: Set[int] = field(default_factory=set)  # 已执行行号集合
    executed_branches: Dict[int, CodeBranch] = field(default_factory=dict)  # 已执行分支
    function_calls: List[FunctionCall] = field(default_factory=list)  # 函数调用列表

# 数据类: 覆盖率报告
# 汇总所有文件的覆盖率统计信息
@dataclass
class CoverageReport:
    total_branches: int = 0                           # 总分支数
    executed_branches: int = 0                        # 已执行分支数
    branches_by_type: Dict[BranchType, int] = field(default_factory=dict)  # 按类型统计
    coverage_percentage: float = 0.0                  # 覆盖率百分比
    file_reports: Dict[str, ExecutionPath] = field(default_factory=dict)  # 各文件报告
\end{lstlisting}

\section{运行时追踪器 (tracer.py)}
\label{sec:tracer}

\begin{lstlisting}[language=Python, caption={运行时追踪器实现 (pathtracer/tracer.py)}]
"""运行时追踪器: 利用 Python 的 sys.settrace 机制追踪程序执行路径"""

import sys
from typing import Optional, Callable, Set, Dict, List
from .models import ExecutionPath, FunctionCall

class PathTracer:
    """基于 sys.settrace 的 Python 程序执行路径追踪器"""

    def __init__(self):
        self.execution_paths: Dict[str, ExecutionPath] = {}  # 各文件的执行路径
        self._original_trace: Optional[Callable] = None      # 原始追踪函数 (用于恢复)
        self._target_file: Optional[str] = None              # 目标追踪文件
        self.function_calls: List[FunctionCall] = []         # 顶层函数调用列表
        self._call_stack: List[FunctionCall] = []            # 函数调用栈 (用于追踪嵌套)

    def trace_file(self, file_path: str) -> 'PathTracer':
        """设置要追踪的目标文件, 返回 self 支持链式调用"""
        self._target_file = file_path
        self.execution_paths[file_path] = ExecutionPath(file_path=file_path)
        return self

    def _trace_callback(self, frame, event: str, arg) -> Optional[Callable]:
        """sys.settrace 的回调函数, 处理各类执行事件"""
        if not self._target_file:
            return None

        # 只追踪目标文件中的事件
        code = frame.f_code
        if not code.co_filename.endswith(self._target_file):
            return self._trace_callback

        # 处理 'line' 事件: 记录执行的代码行
        if event == 'line':
            line_no = frame.f_lineno
            self.execution_paths[self._target_file].executed_lines.add(line_no)

        # 处理 'call' 事件: 记录函数调用信息
        elif event == 'call':
            func_name = code.co_name
            line_no = frame.f_lineno

            # 从栈帧中提取函数参数
            args = {}
            try:
                arg_count = code.co_argcount
                arg_names = code.co_varnames[:arg_count]
                for arg_name in arg_names:
                    if arg_name in frame.f_locals:
                        try:
                            args[arg_name] = repr(frame.f_locals[arg_name])
                        except Exception:
                            args[arg_name] = '<unable to repr>'
            except Exception:
                pass

            # 创建 FunctionCall 对象
            call_depth = len(self._call_stack)
            func_call = FunctionCall(
                function_name=func_name,
                file_path=code.co_filename,
                line_no=line_no,
                call_depth=call_depth,
                arguments=args
            )

            # 将调用添加到父调用的 children 或顶层列表
            if self._call_stack:
                self._call_stack[-1].children.append(func_call)
            else:
                self.function_calls.append(func_call)
                self.execution_paths[self._target_file].function_calls.append(func_call)

            # 压入调用栈
            self._call_stack.append(func_call)

        # 处理 'return' 事件: 记录函数返回值
        elif event == 'return':
            if self._call_stack:
                func_call = self._call_stack.pop()
                if arg is not None:
                    try:
                        func_call.return_value = repr(arg)
                    except Exception:
                        func_call.return_value = '<unable to repr>'

        return self._trace_callback

    def __enter__(self):
        """上下文管理器入口: 启动追踪"""
        self._original_trace = sys.gettrace()
        sys.settrace(self._trace_callback)
        return self

    def __exit__(self, exc_type, exc_val, exc_tb):
        """上下文管理器出口: 停止追踪"""
        sys.settrace(self._original_trace)
        return False

    def get_executed_lines(self, file_path: str) -> Set[int]:
        """获取指定文件的已执行行号集合"""
        if file_path in self.execution_paths:
            return self.execution_paths[file_path].executed_lines
        return set()

    def get_function_calls(self) -> List[FunctionCall]:
        """获取记录的所有函数调用"""
        return self.function_calls
\end{lstlisting}

\section{控制流图构建器 (cfg\_builder.py)}
\label{sec:cfg-builder}

\begin{lstlisting}[language=Python, caption={控制流图构建器实现 (pathtracer/cfg\_builder.py)}]
"""控制流图构建器: 通过静态分析 AST 提取代码中的所有分支点"""

import ast
from typing import List, Optional
from .models import CodeBranch, BranchType

class CFGBuilder:
    """从 Python 源代码构建控制流图, 识别所有分支结构"""

    def __init__(self):
        self.branches: List[CodeBranch] = []  # 提取的分支列表

    def build_from_file(self, file_path: str) -> List[CodeBranch]:
        """解析 Python 文件并提取所有分支点"""
        with open(file_path, 'r', encoding='utf-8') as f:
            source = f.read()

        # 解析源代码为 AST
        tree = ast.parse(source, filename=file_path)
        self.branches = []

        # 遍历 AST 查找控制流结构
        for node in ast.walk(tree):
            branch = self._extract_branch(file_path, node)
            if branch:
                self.branches.append(branch)

        return self.branches

    def _extract_branch(self, file_path: str, node: ast.AST) -> Optional[CodeBranch]:
        """从 AST 节点提取分支信息"""

        # if 语句
        if isinstance(node, ast.If):
            return CodeBranch(
                file_path=file_path,
                line_no=node.lineno,
                branch_type=BranchType.IF,
                end_line=self._get_end_line(node)
            )

        # for 循环
        elif isinstance(node, ast.For):
            return CodeBranch(
                file_path=file_path,
                line_no=node.lineno,
                branch_type=BranchType.FOR,
                end_line=self._get_end_line(node)
            )

        # while 循环
        elif isinstance(node, ast.While):
            return CodeBranch(
                file_path=file_path,
                line_no=node.lineno,
                branch_type=BranchType.WHILE,
                end_line=self._get_end_line(node)
            )

        # try/except 异常处理
        elif isinstance(node, ast.ExceptHandler):
            return CodeBranch(
                file_path=file_path,
                line_no=node.lineno,
                branch_type=BranchType.EXCEPT,
                end_line=self._get_end_line(node)
            )

        return None

    def _get_end_line(self, node: ast.AST) -> int:
        """获取代码块的最后一行行号"""
        # Python 3.8+ 支持 end_lineno 属性
        if hasattr(node, 'end_lineno') and node.end_lineno is not None:
            return node.end_lineno
        # 递归查找 body 中最大的行号
        if hasattr(node, 'body') and node.body:
            return max(self._get_end_line(n) for n in node.body)
        return getattr(node, 'lineno', 0)

    def get_branches_by_type(self, branch_type: BranchType) -> List[CodeBranch]:
        """按类型筛选分支"""
        return [b for b in self.branches if b.branch_type == branch_type]
\end{lstlisting}

\section{报告生成器 (reporter.py)}
\label{sec:reporter}

\begin{lstlisting}[language=Python, caption={报告生成器实现 (pathtracer/reporter.py)}]
"""覆盖率报告生成器: 对比运行时追踪结果与静态 CFG 分析, 生成覆盖率报告"""

import json
from typing import Dict, Any, List
from .models import CodeBranch, ExecutionPath, CoverageReport, BranchType, FunctionCall
from .cfg_builder import CFGBuilder
from .tracer import PathTracer

class CoverageReporter:
    """生成覆盖率报告: 比较追踪路径与控制流图"""

    def __init__(self):
        self.cfg_builder = CFGBuilder()  # 用于静态分析

    def analyze(self, file_path: str, tracer: PathTracer) -> CoverageReport:
        """分析指定文件的覆盖率"""
        # 静态分析: 获取所有分支
        all_branches = self.cfg_builder.build_from_file(file_path)

        # 运行时数据: 获取已执行的行
        executed_lines = tracer.get_executed_lines(file_path)

        # 标记已执行的分支
        for branch in all_branches:
            if branch.line_no in executed_lines:
                branch.is_executed = True

        # 计算统计数据
        total = len(all_branches)
        executed = sum(1 for b in all_branches if b.is_executed)

        # 按分支类型统计
        by_type: Dict[BranchType, int] = {}
        for branch_type in BranchType:
            type_branches = [b for b in all_branches if b.branch_type == branch_type]
            by_type[branch_type] = len(type_branches)

        # 构建并返回覆盖率报告
        return CoverageReport(
            total_branches=total,
            executed_branches=executed,
            branches_by_type=by_type,
            coverage_percentage=(executed / total * 100) if total > 0 else 0.0,
            file_reports={file_path: ExecutionPath(
                file_path=file_path,
                executed_lines=executed_lines,
                executed_branches={b.line_no: b for b in all_branches if b.is_executed}
            )}
        )

    def format_report(self, report: CoverageReport) -> str:
        """将覆盖率报告格式化为可读文本"""
        lines = [
            "=" * 60,
            "EXECUTION PATH COVERAGE REPORT",
            "=" * 60,
            "",
            f"Total Branches: {report.total_branches}",
            f"Executed Branches: {report.executed_branches}",
            f"Coverage: {report.coverage_percentage:.1f}%",
            "",
            "Branches by Type:",
        ]

        # 按类型输出分支统计
        for branch_type, count in report.branches_by_type.items():
            if count > 0:
                lines.append(f"  {branch_type.value:12s}: {count}")

        lines.extend([
            "",
            "=" * 60,
            "EXECUTED BRANCHES:",
            "-" * 60,
        ])

        # 输出各文件的已执行分支详情
        for file_path, exec_path in report.file_reports.items():
            lines.append(f"\n{file_path}:")
            for line_no, branch in sorted(exec_path.executed_branches.items()):
                lines.append(
                    f"  Line {branch.line_no:4d}: {branch.branch_type.value:10s} [executed]"
                )

        return "\n".join(lines)

    def format_report_json(self, report: CoverageReport) -> str:
        """将覆盖率报告格式化为 JSON (便于程序处理)"""
        data: Dict[str, Any] = {
            "coverage_percentage": report.coverage_percentage,
            "total_branches": report.total_branches,
            "executed_branches": report.executed_branches,
            "branches_by_type": {
                bt.value: count for bt, count in report.branches_by_type.items()
            },
            "file_reports": {},
            "function_calls": []
        }

        # 构建各文件的报告数据
        for file_path, exec_path in report.file_reports.items():
            file_data: Dict[str, Any] = {
                "executed_lines": sorted(list(exec_path.executed_lines)),
                "executed_branches": []
            }
            for line_no, branch in sorted(exec_path.executed_branches.items()):
                file_data["executed_branches"].append({
                    "line_no": branch.line_no,
                    "branch_type": branch.branch_type.value,
                    "end_line": branch.end_line
                })
            data["file_reports"][file_path] = file_data

            # 包含函数调用信息
            if exec_path.function_calls:
                data["function_calls"] = self._serialize_function_calls(
                    exec_path.function_calls
                )

        return json.dumps(data, indent=2)

    def _serialize_function_calls(self, calls: List[FunctionCall]) -> List[Dict[str, Any]]:
        """递归序列化函数调用树为 JSON 可用的字典列表"""
        result = []
        for call in calls:
            call_dict = {
                "function_name": call.function_name,
                "file_path": call.file_path,
                "line_no": call.line_no,
                "call_depth": call.call_depth,
                "arguments": call.arguments,
                "return_value": call.return_value,
                "children": self._serialize_function_calls(call.children)  # 递归处理子调用
            }
            result.append(call_dict)
        return result
\end{lstlisting}

\section{命令行接口 (cli.py)}
\label{sec:cli}

\begin{lstlisting}[language=Python, caption={命令行接口实现 (pathtracer/cli.py)}]
"""PathTracer 命令行接口: 提供用户友好的命令行交互方式"""

import sys
import argparse
import io
from pathlib import Path
from .tracer import PathTracer
from .reporter import CoverageReporter

def main():
    """CLI 主入口函数"""
    raw_argv = sys.argv[1:]

    # 无参数时显示帮助信息
    if not raw_argv:
        parser = argparse.ArgumentParser(
            description='Trace Python program execution paths and generate coverage reports'
        )
        parser.add_argument('program', help='Python program file to trace')
        parser.add_argument('--output', '-o', help='Output file for report', default=None)
        parser.add_argument('--quiet', '-q', help='Only output report', action='store_true')
        parser.add_argument(
            '--format', '-f',
            choices=['text', 'html', 'json'],
            default='text',
            help='Output format: text, html, or json (default: text)'
        )
        parser.add_argument('program_args', nargs='*', help='Arguments for target program')
        parser.parse_args(raw_argv)
        return

    # 解析 CLI 参数与目标程序参数 (以第一个位置参数为分界)
    cli_argv = []
    program_args = []
    program = None
    index = 0
    while index < len(raw_argv):
        arg = raw_argv[index]
        # 处理带值的选项
        if arg in ('--output', '-o', '--format', '-f'):
            if index + 1 >= len(raw_argv):
                cli_argv.append(arg)
                break
            cli_argv.extend([arg, raw_argv[index + 1]])
            index += 2
            continue
        # 处理标志选项
        if arg in ('--quiet', '-q'):
            cli_argv.append(arg)
            index += 1
            continue
        # 第个位置参数是目标程序, 后续为其参数
        program = arg
        cli_argv.append(arg)
        program_args = raw_argv[index + 1:]
        break

    # 构建参数解析器
    parser = argparse.ArgumentParser(
        description='Trace Python program execution paths and generate coverage reports'
    )
    parser.add_argument('program', help='Python program file to trace')
    parser.add_argument('--output', '-o', help='Output file for report', default=None)
    parser.add_argument('--quiet', '-q', help='Only output report', action='store_true')
    parser.add_argument(
        '--format', '-f',
        choices=['text', 'html', 'json'],
        default='text',
        help='Output format: text, html, or json (default: text)'
    )
    parser.add_argument('program_args', nargs='*', help='Arguments for target program')

    args = parser.parse_args(cli_argv)

    # 检查目标文件是否存在
    if not Path(args.program).exists():
        print(f"Error: File not found: {args.program}", file=sys.stderr)
        sys.exit(1)

    # 处理参数分隔符 --
    if program_args[:1] == ['--']:
        program_args = program_args[1:]

    # 创建追踪器并设置目标文件
    tracer = PathTracer().trace_file(args.program)

    # 静默模式: 捕获目标程序的 stdout
    if args.quiet:
        old_stdout = sys.stdout
        sys.stdout = io.StringIO()

    old_argv = sys.argv
    reporter = CoverageReporter()

    try:
        # 设置 sys.argv 以便目标程序获取正确参数
        sys.argv = [args.program] + program_args
        # 在追踪上下文中执行目标程序
        with tracer:
            with open(args.program) as f:
                code = compile(f.read(), args.program, 'exec')
                exec(code, {'__name__': '__main__', '__file__': args.program})
    finally:
        # 恢复 sys.argv 和 stdout
        sys.argv = old_argv
        if args.quiet:
            sys.stdout = old_stdout

        # 生成覆盖率报告
        report = reporter.analyze(args.program, tracer)

        # 根据格式选择输出方法
        if args.format == 'json':
            output = reporter.format_report_json(report)
        elif args.format == 'html':
            output = reporter.format_report_html(report)
        else:
            output = reporter.format_report(report)

        # 确定输出文件
        output_file = args.output
        if not output_file:
            if args.format == 'html':
                output_file = 'coverage-report.html'
            elif args.format == 'json':
                output_file = 'coverage-report.json'

        # 输出报告
        if output_file:
            with open(output_file, 'w') as f:
                f.write(output)
            print(f"Report written to: {output_file}")
        else:
            print(output)

if __name__ == '__main__':
    main()
\end{lstlisting}
